\documentclass[a4paper,11pt]{article}

\usepackage[margin=0.2in]{geometry}

\renewcommand{\baselinestretch}{1.0}
\pagenumbering{gobble}
\setlength{\parindent}{0pt}
\setlength{\itemsep}{0.1pt}

\begin{document}

\begin{center}
{\LARGE \textbf{CURRICULUM VITAE}}\\
\vspace{4pt}
{{applicant}}\\
{{home}}\\
No HP: {{phone}} \quad Email: {{email}}
\end{center}

\vspace{8pt}

\textbf{RINGKASAN PROFIL}

Industrial Engineering di industri garment ekspor dengan pengalaman dalam time study, line balancing, capacity planning, dan improvement produksi berbasis data. Terbiasa menangani target output harian, efisiensi line, serta analisa bottleneck untuk peningkatan produktivitas.

\vspace{8pt}

\textbf{KEAHLIAN TEKNIS}

\begin{itemize}
\item Time Study dan Work Measurement
\item Line Balancing dan Efficiency Improvement
\item Capacity Planning dan Output Analysis
\item Bottleneck Identification dan Process Improvement
\item Standard Time Calculation
\item Microsoft Excel (Pivot, Formula, Data Analysis)
\item Production Monitoring (Daily Output, Efficiency)
\item Koordinasi PPC, Produksi, QC
\end{itemize}

\vspace{8pt}

\textbf{PENGALAMAN KERJA}

\textbf{Management Trainee} \\
PT. Eratex Djaja Tbk. Probolinggo \hfill Tahun 2018 - Tahun 2020 \\

\textbf{Industrial Engineering Staff} \\
PT. Eratex Djaja Tbk. Probolinggo \hfill Tahun 2020 - Tahun 2025

\begin{itemize}
\item Perencanaan alur kerja dan permesinan
\item Melakukan time study untuk penentuan standard time operator dan proses produksi
\item Melakukan line balancing untuk optimasi output dan efisiensi lini produksi
\item Analisa bottleneck produksi dan implementasi improvement workflow
\item Monitoring efficiency line, output harian, dan kapasitas produksi
\item Support PPC dalam perencanaan produksi berbasis data aktual
\item Koordinasi dengan produksi dan QC untuk pencapaian target shipment
\item Mengendalikan, mengawasi dan troubleshooting sistem produksi berbasis computerized hanger conveying  system
\end{itemize}

\vspace{8pt}

\textbf{PENCAPAIAN}

\begin{itemize}
\item Peningkatan stabilitas efisiensi line produksi melalui perbaikan proses kerja dan manpower balancing
\item Pengurangan idle time pada proses produksi melalui identifikasi bottleneck
\item Digitalisasi Laporan
\item Terlibat dalam proyek penyusunan Skill Matrix
\item Terlibat dalam proyek Kaizen (Continuous Improvement)
\item Terlibat dalam proyek penerapan tracking system berbasis QRCode\\
\end{itemize}
\vspace{8pt}

\textbf{BAHASA}

\begin{itemize}
\item Bahasa Indonesia (Aktif)
\item Bahasa Inggris (Reading, Listening, Teknis)
\end{itemize}

\end{document}
